% Cover letter for the JSAMS submission. The prose is kept in step with
% manuscript/cover_letter.md; the counts it states are enforced against
% manuscript.tex by tests/test_manuscript_tier_visibility.py.
\documentclass[11pt]{article}

\usepackage[margin=1in]{geometry}
\usepackage{newtxtext}
\usepackage{microtype}
\PassOptionsToPackage{hyphens}{url}
\usepackage[hidelinks]{hyperref}

\setlength{\parindent}{0pt}
\setlength{\parskip}{0.7em}

\begin{document}

To the Editor-in-Chief, \emph{Journal of Science and Medicine in Sport}

Dear Editor,

Please consider our manuscript, \emph{``Dividing by playing time removes part
of the association it should scale: measuring the denominator gradient in
eight European leagues''}, for publication as Original Research.

Category: Original Research. Sub-discipline: Sport Injury (injury
epidemiology). This research did not receive any specific grant from funding
agencies in the public, commercial, or not-for-profit sectors, and there is no
external financial support, equipment or product involvement to declare. The
manuscript is available on the arXiv preprint server as arXiv:2608.11228
(\url{https://arxiv.org/abs/2608.11228}). In line with the journal's preprint
policy, a note will be added to the preprint recording that the manuscript is
under peer review once it is sent for review, and updated again on any
publication decision.

\textbf{Why this journal}

This paper completes an argument this journal began. Shrier and colleagues
showed analytically in the \emph{Journal of Science and Medicine in Sport}
(2022) that when a risk factor changes workload, entering workload ---
including time at risk --- in the regression causes bias. Hoenig and
colleagues, also in this journal (2022), established that the numerator of
public football injury data cannot carry epidemiological weight. What neither
paper could say is how large the denominator's version of the problem is,
where it bites, and how a reader would detect it before dividing. We measure
it: the denominator gradient --- the slope of log recorded minutes on the
exposure --- separates outcome truncation (vivid, and negligible in aggregate)
from exposure-dependent playing time (invisible per row, and decisive), and
replicates without exception across 628,487 appearances in eight European
leagues. The check requires appearance records only --- no injury data --- so
it can be run in any league, before any per-minute rate is computed. We
believe the journal that published the analytical claim is the right home for
its empirical measurement.

\textbf{Ethics and data protection}

The study reused exclusively public, pre-existing records, with no
recruitment, contact, intervention or private medical information. Trinity
College Dublin guidance does not permit retrospective approval of secondary
public-data research; no determination was obtained before the analysis began
and none can now be issued, and the manuscript states this plainly rather than
implying an approval that does not exist. Because the records concern
identifiable professional athletes and absence labels are health-related, the
manuscript also states the data-protection basis: processing for scientific
research under the research provisions of the General Data Protection
Regulation, restricted to data already published by the source.

No deposited file carries a player's name, the data provider's player
identifier, or a match date. The hand-built audit records are keyed by
surrogates drawn at random once and held apart from the archive, with dates
coarsened to seasons; the reviewer's identified copy stays with the authors,
and row-level model diagnostics are not deposited. We do not describe this as
anonymisation, because the sampling rule is published and a reader holding the
data provider's own snapshot could redraw the same appearances. We are happy
to provide any correspondence with our institution the editorial office may
require, and to discuss the de-identification with the editorial office if
that would help.

\textbf{Reproducibility}

Every number in the manuscript and supplement is regenerated from committed
data by archived code, enforced by an automated test suite at 100\% statement
and branch coverage that fails whenever text and data disagree. The frozen
code, data and manuscript are archived together as a single record on Zenodo
and cited in the data-availability statement.

The manuscript is within the journal's 5,000-word limit, with a 249-word
structured abstract, five Practical Implications bullets, two figures and
three tables; the exact counts are stated on the title page and are enforced
against the manuscript by the archived test suite. It is not under
consideration elsewhere, and all authors have approved the submission. The
anonymised, line-numbered review copy, the anonymised supplement and the
separate title page are uploaded per the journal's anonymised review process;
the anonymised files carry neither author names nor the institution's.

Yours sincerely,

Gustavo Pedro Ricou (corresponding author), on behalf of both authors\\
School of Computer Science and Statistics, Trinity College Dublin\\
\href{mailto:pedrorig@tcd.ie}{pedrorig@tcd.ie}

\end{document}
